\begin{ex}
Số liền sau của số 99 là số nào?
\shortans{100}
\loigiai{
Số liền sau của 99 là 100.
}
\end{ex}

\begin{ex}
Số liền trước của số 100 là số nào?
\shortans{99}
\loigiai{
Số liền trước của 100 là 99.
}
\end{ex}

\begin{ex}
Số gồm 7 chục và 4 đơn vị là số nào?
\shortans{74}
\loigiai{
Số gồm 7 chục và 4 đơn vị là 74.
}
\end{ex}

\begin{ex}
Số 85 có chữ số hàng chục là bao nhiêu?
\shortans{8}
\loigiai{
Chữ số hàng chục của 85 là 8.
}
\end{ex}

\begin{ex}
Số 62 có chữ số hàng đơn vị là bao nhiêu?
\shortans{2}
\loigiai{
Chữ số hàng đơn vị của 62 là 2.
}
\end{ex}

\begin{ex}
Số lớn nhất có hai chữ số là số nào?
\shortans{99}
\loigiai{
Số lớn nhất có hai chữ số là 99.
}
\end{ex}

\begin{ex}
Số bé nhất có hai chữ số là số nào?
\shortans{10}
\loigiai{
Số bé nhất có hai chữ số là 10.
}
\end{ex}

\begin{ex}
Trong các số $45, 83, 38, 79$, số nào lớn nhất?
\shortans{83}
\loigiai{
Số lớn nhất là 83.
}
\end{ex}

\begin{ex}
Trong các số $62, 59, 70, 54$, số nào bé nhất?
\shortans{54}
\loigiai{
Số bé nhất là 54.
}
\end{ex}

\begin{ex}
Điền số thích hợp: $59 < \dots < 61$.
\shortans{60}
\loigiai{
Số tự nhiên ở giữa 59 và 61 là 60.
}
\end{ex}

\begin{ex}
Có bao nhiêu số tròn chục có hai chữ số?
\shortans{9}
\loigiai{
Các số tròn chục từ 10 đến 90 gồm 9 số.
}
\end{ex}

\begin{ex}
Có bao nhiêu số có hai chữ số giống nhau?
\shortans{9}
\loigiai{
Gồm 9 số: 11, 22, 33, 44, 55, 66, 77, 88, 99.
}
\end{ex}

\begin{ex}
Sắp xếp các số $35, 82, 19, 64$ từ bé đến lớn, số đứng ở vị trí thứ hai là số nào?
\shortans{35}
\loigiai{
Thứ tự: 19, 35, 64, 82. Số thứ hai là 35.
}
\end{ex}

\begin{ex}
Sắp xếp các số $71, 48, 93, 65$ từ lớn đến bé, số đứng ở vị trí thứ ba là số nào?
\shortans{65}
\loigiai{
Thứ tự: 93, 71, 65, 48. Số thứ ba là 65.
}
\end{ex}

\begin{ex}
Từ 20 đến 30 có tất cả bao nhiêu số tự nhiên?
\shortans{11}
\loigiai{
Có $(30 - 20) + 1 = 11$ số.
}
\end{ex}

\begin{ex}
Tìm số lớn nhất có hai chữ số mà tổng hai chữ số bằng 9.
\shortans{90}
\loigiai{
Số lớn nhất là 90 vì $9 + 0 = 9$ và có hàng chục lớn nhất.
}
\end{ex}

\begin{ex}
Tìm số bé nhất có hai chữ số mà tổng hai chữ số bằng 9.
\shortans{18}
\loigiai{
Số bé nhất có hàng chục là 1, hàng đơn vị là 8 ($1 + 8 = 9$), vậy là 18.
}
\end{ex}

\begin{ex}
Tính tổng chữ số hàng chục và hàng đơn vị của số 47.
\shortans{11}
\loigiai{
Tổng là: $4 + 7 = 11$.
}
\end{ex}

\begin{ex}
Hiệu giữa chữ số hàng chục và hàng đơn vị của số 82 là bao nhiêu?
\shortans{6}
\loigiai{
Hiệu là: $8 - 2 = 6$.
}
\end{ex}

\begin{ex}
Lớp 2A có 35 học sinh, lớp 2B có 38 học sinh. Lớp 2B nhiều hơn lớp 2A bao nhiêu học sinh?
\shortans{3}
\loigiai{
Nhiều hơn: $38 - 35 = 3$ học sinh.
}
\end{ex}

\begin{ex}
Có bao nhiêu số có một chữ số?
\shortans{10}
\loigiai{
Từ 0 đến 9 gồm 10 số.
}
\end{ex}

\begin{ex}
Có bao nhiêu số có hai chữ số?
\shortans{90}
\loigiai{
Từ 10 đến 99 gồm 90 số.
}
\end{ex}

\begin{ex}
Số tròn chục liền sau số 50 là số nào?
\shortans{60}
\loigiai{
Số tròn chục liền sau 50 là 60.
}
\end{ex}

\begin{ex}
Số tròn chục liền trước số 80 là số nào?
\shortans{70}
\loigiai{
Số tròn chục liền trước 80 là 70.
}
\end{ex}

\begin{ex}
Số nào trừ đi 1 thì được số lớn nhất có một chữ số?
\shortans{10}
\loigiai{
Số lớn nhất có 1 chữ số là 9. Số cần tìm: $9 + 1 = 10$.
}
\end{ex}
